\documentclass{amsart}
\usepackage{amsmath, amssymb, amsthm, hyperref, tikz}

\title{Toric Geometry and Calabi--Yau Manifolds: Lecture Notes}
\author{Personalized Notes (companion to \texttt{refpoly-5d})}
\date{\today}

\newtheorem{theorem}{Theorem}[section]
\newtheorem{lemma}[theorem]{Lemma}
\newtheorem{proposition}[theorem]{Proposition}
\newtheorem{corollary}[theorem]{Corollary}
\newtheorem{definition}[theorem]{Definition}
\newtheorem{example}[theorem]{Example}
\newtheorem{remark}[theorem]{Remark}

\begin{document}
\maketitle
\tableofcontents

%==============================================================
\section{Foundations of Toric Geometry}
%==============================================================
We fix throughout a rank-$n$ lattice $N \cong \mathbb{Z}^n$ with dual lattice
$M = \operatorname{Hom}(N,\mathbb{Z})$, and their real scalar extensions
$N_{\mathbb{R}} = N \otimes \mathbb{R}$, $M_{\mathbb{R}} = M \otimes \mathbb{R}$.
The natural pairing is denoted $\langle\,\cdot\,,\,\cdot\,\rangle : M_{\mathbb{R}} \times N_{\mathbb{R}} \to \mathbb{R}$.

A \emph{fan} $\Sigma$ in $N_{\mathbb{R}}$ is a collection of strongly convex
rational polyhedral cones, closed under taking faces and intersections.
To each fan one associates a normal toric variety $\mathbb{P}_\Sigma$
(equivalently $X_\Sigma$). The construction translates combinatorial
data of $\Sigma$ into algebro-geometric data of $X_\Sigma$:
cones $\leftrightarrow$ affine charts, rays $\leftrightarrow$ torus-invariant
divisors, smoothness $\leftrightarrow$ each cone spanned by part of a
lattice basis, compactness $\leftrightarrow$ the support of $\Sigma$ being
all of $N_{\mathbb{R}}$.

\begin{remark}
This section is a placeholder to be filled in as the foundational
toric material (cones, fans, divisors, the Cox ring, intersection theory)
is discussed. The notes below currently focus on Calabi--Yau manifolds,
reflexive polytopes, and Batyrev's construction.
\end{remark}

%==============================================================
\section{Calabi--Yau Manifolds}
%==============================================================

\subsection{Definition and equivalences}

\begin{definition}[Calabi--Yau manifold]\label{def:cy}
A \emph{Calabi--Yau manifold} of complex dimension $n$ is a compact
K\"ahler manifold $X$ satisfying any one of the following equivalent
conditions:
\begin{enumerate}
  \item the canonical bundle is trivial,
        $K_X := \bigwedge^n \Omega_X^1 \cong \mathcal{O}_X$
        (equivalently $c_1(X) = 0$);
  \item the holonomy group of $X$ is contained in $SU(n)$;
  \item there exists a nowhere-vanishing holomorphic $n$-form
        $\Omega \in H^0(X,\Omega_X^n)$, unique up to scale;
  \item each K\"ahler class contains a Ricci-flat K\"ahler metric.
\end{enumerate}
\end{definition}

\begin{theorem}[Yau, 1978; Calabi conjecture]\label{thm:yau}
On a compact K\"ahler manifold with $c_1(X)=0$, every K\"ahler class
contains a \emph{unique} Ricci-flat K\"ahler metric.
\end{theorem}

\begin{remark}
Theorem~\ref{thm:yau} is the bridge between the algebraic condition
$K_X \cong \mathcal{O}_X$ (Definition~\ref{def:cy}(1), checkable by pure
algebra) and the differential-geometric condition (Definition~\ref{def:cy}(4)).
For no compact Calabi--Yau of dimension $\geq 2$ is the Ricci-flat metric
known explicitly; constructing it numerically is an active area.
\end{remark}

\begin{remark}[A definitional subtlety]
Some authors strengthen Definition~\ref{def:cy} by also requiring
$h^{i,0}(X)=0$ for $0<i<n$, which excludes complex tori and products and
forces holonomy to be \emph{exactly} $SU(n)$. This matters when counting
moduli but rarely affects the toric hypersurface constructions below.
\end{remark}

\subsection{Hodge numbers and moduli}

The topology of $X$ is organized by the Hodge numbers
$h^{p,q} = \dim_{\mathbb{C}} H^q(X,\Omega_X^p)$. For a Calabi--Yau
threefold satisfying the strict definition, the Hodge diamond has only two
independent entries, $h^{1,1}$ and $h^{2,1}$, with Euler characteristic
\begin{equation}
  \chi(X) = 2\bigl(h^{1,1} - h^{2,1}\bigr).
\end{equation}
Geometrically and physically:
\begin{itemize}
  \item $h^{1,1}$ counts \emph{K\"ahler (size) moduli} --- deformations of
        the K\"ahler class;
  \item $h^{2,1}$ counts \emph{complex-structure (shape) moduli} ---
        deformations of the complex structure.
\end{itemize}

%==============================================================
\section{Why Calabi--Yau Manifolds Matter}
%==============================================================

\subsection{Mathematical significance}
Calabi--Yau manifolds are the $\kappa = 0$, $K_X \cong \mathcal{O}_X$
fixed point of the canonical class, sitting between Fano varieties
($-K_X$ ample, ``positively curved'') and varieties of general type
($K_X$ ample) in the Kodaira/minimal-model trichotomy. Yau's theorem
(Theorem~\ref{thm:yau}) makes them central to geometric analysis, and
mirror symmetry (Section~\ref{sec:mirror}) makes them a font of
enumerative-geometry predictions.

\subsection{Physical significance}
In superstring theory spacetime is $10$-dimensional; compactification
writes it as $\mathbb{R}^{1,3} \times X$ with $X$ compact and small.
Preserving $\mathcal{N}=1$ supersymmetry in the $4$d effective theory
forces $X$ to have $SU(3)$ holonomy, i.e.\ $X$ is a Calabi--Yau
\emph{threefold} (Candelas--Horowitz--Strominger--Witten, 1985). The
dictionary reads $h^{2,1}, h^{1,1} \leftrightarrow$ massless scalar moduli
fields, and the number of chiral generations is governed by
$\tfrac{1}{2}|\chi(X)|$ in simple heterotic models.

\begin{remark}[Relevance to this repository]\label{rmk:repo}
F-theory compactifications use Calabi--Yau \emph{fourfolds}, which arise
as anticanonical hypersurfaces in $5$-dimensional toric varieties. A
classification of $5$-dimensional reflexive polytopes (the goal of the
\texttt{refpoly-5d} project) is therefore a classification of candidate
ambient spaces for such F-theory geometries.
\end{remark}

%==============================================================
\section{Reflexive Polyhedra}
%==============================================================

\begin{definition}[Reflexive polytope; Batyrev 1994]\label{def:reflexive}
A full-dimensional lattice polytope $\Delta \subset M_{\mathbb{R}}$
containing the origin in its interior is \emph{reflexive} if its
\emph{polar dual}
\begin{equation}
  \Delta^* \;:=\; \bigl\{\, v \in N_{\mathbb{R}} \;:\;
  \langle u,v\rangle \geq -1 \ \ \forall\, u \in \Delta \,\bigr\}
\end{equation}
is again a lattice polytope (all of its vertices lie in $N$).
\end{definition}

\begin{proposition}[Properties of reflexive polytopes]\label{prop:reflexive}
\leavevmode
\begin{enumerate}
  \item Polar duality is an involution: $(\Delta^*)^* = \Delta$. Hence
        reflexive polytopes occur in \emph{dual pairs} $(\Delta,\Delta^*)$,
        with the roles of $M$ and $N$ exchanged.
  \item A reflexive polytope has \emph{exactly one} interior lattice point,
        namely the origin. (This is the practical test for reflexivity.)
  \item In each dimension there are finitely many reflexive polytopes up to
        $GL(n,\mathbb{Z})$ (Lagarias--Ziegler finiteness).
\end{enumerate}
\end{proposition}

\begin{remark}[Kreuzer--Skarke enumeration]
The classification is complete through dimension $4$:
$1$ in dimension $0$, $16$ reflexive polygons in dimension $2$,
$4319$ in dimension $3$, and $473\,800\,776$ in dimension $4$
(Kreuzer--Skarke). Dimension $5$ is \emph{not} fully classified --- the
frontier addressed by the \texttt{refpoly-5d} project (cf.\
Remark~\ref{rmk:repo}).
\end{remark}

\begin{example}[The reflexive simplex of $\mathbb{P}^2$]\label{ex:p2}
The simplex $\Delta \subset M_{\mathbb{R}} = \mathbb{R}^2$ with vertices
$(-1,-1),\ (2,-1),\ (-1,2)$ is reflexive; its only interior lattice point
is the origin. Its polar dual $\Delta^*$ has vertices
$(1,0),\ (0,1),\ (-1,-1)$ and is also reflexive. The associated toric
variety is $\mathbb{P}^2$.
\end{example}

%==============================================================
\section{Batyrev's Construction and Mirror Symmetry}\label{sec:mirror}
%==============================================================

\subsection{From a reflexive polytope to a Gorenstein Fano toric variety}

Given a reflexive $\Delta \subset M_{\mathbb{R}}$, the face fan
(normal fan) of $\Delta^*$ defines a toric variety $\mathbb{P}_\Delta$.

\begin{theorem}[Reflexivity $\Leftrightarrow$ Gorenstein Fano]\label{thm:fano}
$\Delta$ is reflexive if and only if $\mathbb{P}_\Delta$ is a Gorenstein
Fano toric variety, i.e.\ the anticanonical class $-K_{\mathbb{P}_\Delta}$
is an ample Cartier divisor. The lattice points $\Delta \cap M$ index a
basis of monomial sections of $-K_{\mathbb{P}_\Delta}$.
\end{theorem}

\subsection{The Calabi--Yau hypersurface}

\begin{theorem}[Batyrev, 1994]\label{thm:batyrev}
Let $\Delta \subset M_{\mathbb{R}}$ be reflexive of dimension $n$. A generic
member of the anticanonical linear system,
\begin{equation}
  X \;=\; \Bigl\{\, \textstyle\sum_{m \in \Delta \cap M} a_m\, x^{m} = 0 \,\Bigr\}
  \;\subset\; \mathbb{P}_\Delta,
\end{equation}
admits a crepant toric resolution that is a Calabi--Yau variety of
dimension $n-1$.
\end{theorem}

\begin{proof}[Mechanism (adjunction)]
For $X \in |-K_{\mathbb{P}_\Delta}|$ the adjunction formula gives
\begin{equation}
  K_X = \bigl(K_{\mathbb{P}_\Delta} + X\bigr)\big|_X
      = \bigl(K_{\mathbb{P}_\Delta} - K_{\mathbb{P}_\Delta}\bigr)\big|_X
      = \mathcal{O}_X,
\end{equation}
so the canonical bundle is trivial by construction. Reflexivity
(Theorem~\ref{thm:fano}) ensures $-K_{\mathbb{P}_\Delta}$ is Cartier and
base-point-free, and a sufficiently fine fan subdivision yields a crepant
(Calabi--Yau-preserving) resolution of the residual orbifold
singularities.
\end{proof}

\begin{example}[The quintic threefold]\label{ex:quintic}
Taking $n=4$ and $\Delta$ the reflexive simplex with
$\mathbb{P}_\Delta = \mathbb{P}^4$, one has
$-K_{\mathbb{P}^4} = \mathcal{O}(5)$, so $X$ is a generic quintic
hypersurface in $\mathbb{P}^4$: a Calabi--Yau threefold with
$h^{1,1}=1$, $h^{2,1}=101$, $\chi = -200$.
\end{example}

\subsection{Mirror symmetry via polar duality}

The reflexive pair $(\Delta,\Delta^*)$ yields a \emph{mirror pair} of
Calabi--Yau varieties $(X, X^*)$: applying Theorem~\ref{thm:batyrev} to
$\Delta^*$ gives $X^*$. For the threefold case ($n=4$), Batyrev's
combinatorial Hodge-number formulas read
\begin{align}
  h^{1,1}(X) &= \ell(\Delta^*) - (n+1)
     - \sum_{\substack{\text{facets }\Theta^* \prec \Delta^*}} \ell^*(\Theta^*)
     + \sum_{\substack{\text{codim-}2\ \Theta^* \prec \Delta^*}}
        \ell^*(\Theta^*)\,\ell^*(\widehat{\Theta}),
\end{align}
where $\ell(\cdot)$ counts lattice points, $\ell^*(\cdot)$ counts
relative-interior lattice points, and $\widehat{\Theta}$ is the dual face;
the formula for $h^{2,1}(X)$ is the same expression with $\Delta$ and
$\Delta^*$ interchanged. The manifest $\Delta \leftrightarrow \Delta^*$
symmetry yields
\begin{equation}
  h^{1,1}(X) = h^{2,1}(X^*), \qquad h^{2,1}(X) = h^{1,1}(X^*),
\end{equation}
the topological signature of mirror symmetry (a reflection of the Hodge
diamond across its diagonal).

\begin{remark}
Batyrev's construction was the first systematic, large-scale source of
mirror pairs, and it is why the Kreuzer--Skarke database is the canonical
arena for studying mirror symmetry. The deeper physical statement --- that
the quantum-corrected K\"ahler moduli space of $X$ equals the
complex-structure moduli space of $X^*$, with Gromov--Witten invariants of
$X$ extracted from period integrals on $X^*$ --- was verified explicitly
for the quintic by Candelas--de la Ossa--Green--Parkes.
\end{remark}

%==============================================================
\section{Key Results from the Literature}
%==============================================================
\begin{itemize}
  \item \textbf{Batyrev (1994)}, ``Dual polyhedra and mirror symmetry for
        Calabi--Yau hypersurfaces in toric varieties,'' \emph{J.\ Algebraic
        Geom.}\ \textbf{3}, 493--535 (arXiv:alg-geom/9310003): reflexive
        polytopes, Theorems~\ref{thm:fano}, \ref{thm:batyrev}, and the
        mirror construction.
  \item \textbf{Yau (1978)}: existence and uniqueness of Ricci-flat K\"ahler
        metrics (Theorem~\ref{thm:yau}).
  \item \textbf{Candelas--de la Ossa--Green--Parkes (1991)}: explicit mirror
        symmetry and Gromov--Witten predictions for the quintic.
  \item \textbf{Kreuzer--Skarke}: complete enumeration of reflexive
        polytopes through dimension $4$.
\end{itemize}

\end{document}
